% ---------------- Packages ----------------
%\usepackage{hyphenat}

\usepackage[utf8]{inputenc}
\usepackage[T1]{fontenc}
\usepackage{lmodern}
\usepackage{graphicx}
% \usepackage[fleqn]{amsmath}
\usepackage{amsmath}
\usepackage{amsthm}
\usepackage{mdframed}
\usepackage{listings}
\usepackage{booktabs}
\usepackage{hyperref}
% \usepackage{dtsyntax}
% \usepackage[colorinlistoftodos,prependcaption,textsize=tiny]{todonotes}
\usepackage[normalem]{ulem}
\usepackage{pdflscape}
\usepackage{amssymb}
\usepackage{longtable}
\usepackage{tikz}
\usetikzlibrary{positioning, arrows.meta, babel, shapes.multipart, decorations.pathreplacing}
%\usepackage{algorithm}
\usepackage[authoryear,round]{natbib}
\usepackage{mathpartir}
\usepackage{multicol}
\usepackage{centernot}
\usepackage{hyperref}
\usepackage{array}
\usepackage{tabularx}
\renewcommand{\chapterautorefname}{Capítulo}
\renewcommand{\sectionautorefname}{Sección}
\renewcommand{\figureautorefname}{Figura}
\renewcommand{\tableautorefname}{Tabla}
\renewcommand{\arraystretch}{1.4}

\usepackage{placeins}

\usepackage[a4paper,top=2.5cm,bottom=3cm,inner=2.5cm,outer=2.5cm]{geometry}

\usepackage{needspace}
\usepackage{pifont}

\newcommand{\cmark}{\textcolor{green!50!black}{\ding{51}}}
\newcommand{\xmark}{\textcolor{red!70!black}{\ding{55}}}

\usepackage{multirow}
% For \includegraphics with valign
\usepackage[export]{adjustbox}

\usepackage{tablefootnote}

% M: Tabla en vez de Cuadro
\usepackage[spanish,es-tabla]{babel}

\usepackage{xargs} % Use more than one optional parameter in a new commands

% Algorithms boxes
\usepackage{algorithm}
\usepackage{algorithmic}

% For \qtyrange and \hertz
\usepackage{siunitx}
\usepackage{textcomp}
\DeclareSIUnit{\pixel}{px}
\DeclareSIUnit{\fps}{fps}

\usepackage{subfig}
% \usepackage{subcaption}

% F* code
\usepackage{listings}
\usepackage{xcolor}
\renewcommand{\lstlistingname}{Código}

\definecolor{fstarnavy}{HTML}{254B6E}
\definecolor{fstargreen}{HTML}{356B4A}
\definecolor{fstargray}{HTML}{666666}

% Replace complete keyword tokens without changing longer identifiers.
\makeatletter
\newcommand{\fstarsymbol}[1]{%
  \color{black}%
  \lst@token{\ensuremath{#1}}%
  \lst@length=1\relax
}
\makeatother

\lstdefinelanguage{FStar}{
  sensitive=true,
  alsoletter={'},
  morekeywords={
    module,open,include,
    let,rec,and,in,
    type,noeq,unfold,
    match,with,when,
    if,then,else,
    function,
    val,assume,
    requires,ensures,decreases,
    Lemma,Tot,GTot,Type,Type0,prop,
    squash
  },
  morekeywords=[2]{exists},
  morekeywords=[3]{forall},
  morekeywords=[4]{fun},
  morecomment=[l]{//},
  morecomment=[s]{(*}{*)},
  morestring=[b]"
}

\lstdefinestyle{fstar}{
  language=FStar,
  literate=*
    {/\\}{{\color{black}\ensuremath{\land}}}1
    {\\/}{{\color{black}\ensuremath{\lor}}}1
    {->}{{\color{black}\ensuremath{\rightarrow}}}1
    {^->}{{\color{black}\ensuremath{\rightsquigarrow}}}1
    {==>}{{\color{black}\ensuremath{\Rightarrow}}}1
    {<==>}{{\color{black}\ensuremath{\Leftrightarrow}}}1,
  basicstyle=\ttfamily\small,
  breakindent=1.5em,
  keywordstyle=\color{fstarnavy},
  keywordstyle=[2]\fstarsymbol{\exists},
  keywordstyle=[3]\fstarsymbol{\forall},
  keywordstyle=[4]\fstarsymbol{\lambda},
  commentstyle=\color{fstargray}\itshape,
  stringstyle=\color{fstargreen},
  %numbers=left,
  %numberstyle=\tiny\color{fstargray},
  %numbersep=10pt,
  %stepnumber=1,
  frame=none,
  breaklines=true,
  breakatwhitespace=true,
  showstringspaces=false,
  keepspaces=true,
  columns=fullflexible,
  tabsize=2,
  xleftmargin=1.8em,
  framexleftmargin=1.8em,
  captionpos=b
}

\lstset{style=fstar}

\mdfdefinestyle{barraizquierda}{
  hidealllines=true,
  leftline=true,
  linecolor=black,
  linewidth=0.8pt,
  innerleftmargin=10pt,
  innerrightmargin=0pt,
  innertopmargin=3pt,
  innerbottommargin=3pt,
  skipabove=10pt,
  skipbelow=10pt
}

% ---------------- Macros / Custom commands ----------------
\newcommand{\fstar}[0]{\texorpdfstring{F$^\star$}{F*}}
\newcommand{\nmapsto}{\centernot \mapsto}

% Title page
\newcommand*{\titleGP}{\begingroup
    \centering
    \vspace*{\baselineskip}
    \thispagestyle{empty}

    \rule{\textwidth}{1.6pt}\vspace*{-\baselineskip}\vspace*{2pt}
    \rule{\textwidth}{0.4pt}\\[\baselineskip]

    {\LARGE Formalización de Lógica de Incorrectitud de Separación en \fstar } \\[0.2\baselineskip]

    \rule{\textwidth}{0.4pt}\vspace*{-\baselineskip}\vspace{3.2pt}
    \rule{\textwidth}{1.6pt}\\[\baselineskip]

    \vspace*{2\baselineskip}

    \includegraphics[width=0.3\textwidth]{img/logo-unr.png}

    \vspace*{2\baselineskip}

    \scshape
    Tesina de Grado\\[\baselineskip]
    Licenciatura en Ciencias de la Computación\par

    \vspace*{2\baselineskip}

    {\Large Autora: María Agustina Torrano \\}
    {Legajo: T-2989/1 }
    \vspace*{2\baselineskip}
    {\Large \\ Director: Guido Martínez } \vfill

    \vspace*{2\baselineskip}
    {\scshape \today} \\[0.3\baselineskip]

\endgroup}

\newcommand{\triple}[3]{ \left\{ #1 \right\}~#2~\left\{ #3 \right\} }
\newcommand{\itriple}[3]{ \left\lbrack #1 \right\rbrack~#2~\left\lbrack #3 \right\rbrack }
\theoremstyle{definition}
\newtheorem{definicion}{Definición}[chapter]
\theoremstyle{plain}
\newtheorem{teorema}[definicion]{Teorema}
\newtheorem{lema}[definicion]{Lema}
\newtheorem{corolario}[definicion]{Corolario}
\newtheorem{ejemplo}{Ejemplo}[chapter]
\newtheorem{observacion}[definicion]{Observación}
\surroundwithmdframed[style=barraizquierda]{definicion}
\surroundwithmdframed[style=barraizquierda]{teorema}
\surroundwithmdframed[style=barraizquierda]{lema}
\surroundwithmdframed[style=barraizquierda]{corolario}
